% Part: first-order-logic
% Chapter: tableaux
% Section: soundness.tex

\documentclass[../../../include/open-logic-section]{subfiles}

\begin{document}

\iftag{FOL}
      {\olfileid{fol}{tab}{sou}}
      {\olfileid{pl}{tab}{sou}}
      
\olsection{निर्दोषता}

\begin{explain}
टॅब्लोसारखी !!^a{derivation} पद्धत \emph{निर्दोष} असते, जर ती प्रत्यक्षात
धारण न करणाऱ्या गोष्टी !!{derive} करू शकत नसेल. त्यामुळे निर्दोषता हा
!!{derivation} पद्धतींसाठी हमी दिलेल्या सुरक्षिततेचा एक प्रकार आहे. कोणता
सिद्धता-उपपत्तीय गुणधर्म विचारात आहे यावर अवलंबून, उदाहरणार्थ, आपल्याला
पुढील गोष्टी हव्या असतात:
\begin{enumerate}
\item प्रत्येक !!{derivable}~$!A$ हे \iftag{FOL}{वैध}{सर्वतः सत्य सूत्र} असते;
\item एखादे !!a{sentence} इतर काही वाक्यांपासून !!{derivable} असेल, तर ते
  त्यांचा तार्किक परिणामही असते;
\item !!{sentence}s चा संच विसंगत असेल, तर तो अपूर्ततायोग्य असतो.
\end{enumerate}
हे !!a{derivation} पद्धतीचे महत्त्वाचे गुणधर्म आहेत. यांपैकी एखादा गुणधर्म
खरा नसेल, तर !!{derivation} पद्धतीत दोष आहे---ती खूप जास्त गोष्टी
!!{derive} करेल. म्हणून !!a{derivation} पद्धतीची निर्दोषता सिद्ध करणे अत्यंत
महत्त्वाचे आहे.

हे सर्व सिद्धता-उपपत्तीय गुणधर्म कोणत्या ना कोणत्या प्रकारच्या बंद
!!{tableau}s द्वारे परिभाषित केलेले असल्यामुळे, वरील (1)--(3) सिद्ध करण्यासाठी
बंद !!{tableau}s च्या चिन्हार्थविषयक गुणधर्मांबद्दल काहीतरी सिद्ध करावे लागते.
प्रथम एखाद्या !!a{signed
  formula} ची संरचनेत पूर्ती होणे म्हणजे काय ते
परिभाषित करू; मग !!{tableau} बंद असेल, तर कोणतीही संरचना त्याची सर्व गृहीतके
पूर्ण करत नाही, हे दाखवू. त्यानंतर (1)--(3) या परिणामाचे उपपरिणाम म्हणून
मिळतील.
\end{explain}

\begin{defn}
\iftag{FOL}{!!^a{structure}}{!!^a{valuation}}~$\iftag{FOL}{\Struct{M}}{\pAssign{v}}$
!!a{signed formula} $\sFmla{\True}{!A}$ \emph{पूर्ण करते} तेव्हा आणि केवळ तेव्हाच, जेव्हा
$\iftag{FOL}{\Sat{M}}{\pSat{v}}{!A}$; आणि ते
$\sFmla{\False}{!A}$ पूर्ण करते तेव्हा आणि केवळ तेव्हाच, जेव्हा
$\iftag{FOL}{\Sat/{M}}{\pSat/{v}}{!A}$. $\iftag{FOL}{\Struct{M}}{\pAssign{v}}$
हा !!{signed formula}s चा संच~$\Gamma$ पूर्ण करतो तेव्हा आणि केवळ तेव्हाच,
जेव्हा तो प्रत्येक $\sFmla{S}{!A} \in \Gamma$ पूर्ण करतो. $\Gamma$ पूर्ण
करणारी एखादी \iftag{FOL}{!!a{structure}}{!!a{valuation}} असेल, तर तो
\emph{पूर्ततायोग्य} आणि अन्यथा \emph{अपूर्ततायोग्य} असतो.
\end{defn}

\begin{thm}[निर्दोषता]
  \ollabel{thm:tableau-soundness}
  $\Gamma$ साठी बंद !!{tableau} असेल, तर $\Gamma$ अपूर्ततायोग्य असतो.
\end{thm}

\begin{proof}
!!a{tableau} च्या शाखेवरील !!{signed formula}s चा संच पूर्ततायोग्य असेल, तर
आणि केवळ तेव्हाच त्या शाखेला पूर्ततायोग्य म्हणू; तसेच !!a{tableau} मध्ये
किमान एक पूर्ततायोग्य शाखा असेल, तर त्याला पूर्ततायोग्य म्हणू.

पुढील गोष्ट दाखवू: एखाद्या पूर्ततायोग्य !!{tableau} चा अनुमानाच्या कोणत्याही
एका नियमाने विस्तार केल्यास मिळणारा !!{tableau} नेहमी पूर्ततायोग्य असतो.
यावरून प्रमेय सिद्ध होईल: $\Gamma$ मधील केवळ गृहीतके असलेल्या !!{tableau} ला
अनुमानाचे नियम लावून कोणताही बंद !!{tableau} मिळतो. म्हणून $\Gamma$
पूर्ततायोग्य असेल, तर त्याच्यासाठीचा कोणताही !!{tableau} पूर्ततायोग्य असेल.
पण बंद !!{tableau} स्पष्टपणे पूर्ततायोग्य नसतो: प्रत्येक शाखेत
$\sFmla{\True}{!A}$ आणि $\sFmla{\False}{!A}$ ही दोन्ही सूत्रे असतात, आणि
कोणतीही संरचना $!A$ ची पूर्ती आणि अपूर्ती या दोन्ही करू शकत नाही.

आपल्याकडे पूर्ततायोग्य !!{tableau}, म्हणजे किमान एक पूर्ततायोग्य शाखा असलेला
!!a{tableau}, आहे असे समजा. अनुमानाचा नियम लावल्यावर एकतर शाखेत
!!{signed formula}s जोडली जातात, किंवा शाखा दोन भागांत विभागली जाते. त्या
नियमाच्या उपयोगाने विस्तारली न गेलेली एखादी पूर्ततायोग्य शाखा !!{tableau} मध्ये
असेल, तर विस्तारित !!{tableau} मध्येही ती पूर्ततायोग्य शाखा राहते; म्हणून
विस्तारित टॅब्लोसुद्धा पूर्ततायोग्य असतो. त्यामुळे पूर्ततायोग्य शाखेलाच नियम
लावला जातो ते प्रकरण पाहणे पुरेसे आहे.

त्या शाखेवरील !!{signed formula}s चा संच~$\Gamma$ असू द्या, आणि ज्या
!!{signed formula} ला नियम लावला जातो ते $\sFmla{S}{!A} \in \Gamma$ असू
द्या. नियम लावल्यावर शाखा विभागली जात नसेल, तर विस्तारित शाखा, म्हणजे
$\Gamma$ आणि नियमाचे निष्कर्ष मिळून बनलेला संच, अजूनही पूर्ततायोग्य आहे हे
दाखवावे लागेल. नियम लावल्यावर शाखा विभागली जात असेल, तर मिळणाऱ्या दोन
शाखांपैकी किमान एक पूर्ततायोग्य आहे हे दाखवावे लागेल.
  
प्रथम, शाखा न विभागणारी संभाव्य अनुमाने पाहू.
\begin{enumerate}
\item $\sFmla{\True}{\lnot !B} \in \Gamma$ ला
  $\TRule{\True}{\lnot}$ लावून शाखेचा विस्तार केला आहे असे समजा. मग
  विस्तारित शाखेवरील !!{signed formula}s चा संच $\Gamma \cup
  \{\sFmla{\False}{!B}\}$ असतो. समजा
  $\iftag{FOL}{\Sat{M}}{\pSat{v}}{\Gamma}$. विशेषतः,
  $\iftag{FOL}{\Sat{M}}{\pSat{v}}{\lnot !B}$. त्यामुळे
  $\iftag{FOL}{\Sat/{M}}{\pSat/{v}}{!B}$, म्हणजे
  $\iftag{FOL}{\Struct{M}}{\pAssign{v}}$ हे
  $\sFmla{\False}{!B}$ पूर्ण करते.
\item $\sFmla{\False}{\lnot !B} \in \Gamma$ ला
  $\TRule{\False}{\lnot}$ लावून शाखेचा विस्तार केला आहे: सराव.
\item $\sFmla{\True}{!B \land !C} \in \Gamma$ ला
  $\TRule{\True}{\land}$ लावून शाखेचा विस्तार केला आहे. यामुळे शाखेवर
  $\sFmla{\True}{!B}$ आणि $\sFmla{\True}{!C}$ ही दोन नवी
  !!{signed formula}s मिळतात. समजा
  $\iftag{FOL}{\Sat{M}}{\pSat{v}}{\Gamma}$; विशेषतः,
  $\iftag{FOL}{\Sat{M}}{\pSat{v}}{!B \land !C}$. मग
  $\iftag{FOL}{\Sat{M}}{\pSat{v}}{!B}$ आणि
  $\iftag{FOL}{\Sat{M}}{\pSat{v}}{!C}$. म्हणजे
  $\iftag{FOL}{\Struct{M}}{\pAssign{v}}$ हे
  $\sFmla{\True}{!B}$ आणि $\sFmla{\True}{!C}$ ही दोन्ही सूत्रे पूर्ण करते.
\item $\sFmla{\False}{!B \lor !C} \in \Gamma$ ला
  $\TRule{\False}{\lor}$ लावून शाखेचा विस्तार केला आहे: सराव.
\item $\sFmla{\False}{!B \lif !C} \in \Gamma$ ला
  $\TRule{\False}{\lif}$ लावून शाखेचा विस्तार केला आहे. यामुळे शाखेवर
  $\sFmla{\True}{!B}$ आणि $\sFmla{\False}{!C}$ ही दोन नवी
  !!{signed formula}s मिळतात. समजा
  $\iftag{FOL}{\Sat{M}}{\pSat{v}}{\Gamma}$; विशेषतः,
  $\iftag{FOL}{\Sat/{M}}{\pSat/{v}}{!B \lif !C}$. मग
  $\iftag{FOL}{\Sat{M}}{\pSat{v}}{!B}$ आणि
  $\iftag{FOL}{\Sat/{M}}{\pSat/{v}}{!C}$. म्हणजे
  $\iftag{FOL}{\Struct{M}}{\pAssign{v}}$ हे
  $\sFmla{\True}{!B}$ आणि $\sFmla{\False}{!C}$ ही दोन्ही सूत्रे पूर्ण करते.

\iftag{FOL}{%
\item $\sFmla{\True}{\lforall[x][!A(x)]} \in \Gamma$ ला
  $\TRule{\True}{\lforall}$ लावून शाखेचा विस्तार केला आहे. यामुळे शाखेवर
  नवीन !!{signed formula}~$\sFmla{\True}{!A(t)}$ मिळते.
  समजा $\Sat{M}{\Gamma}$; विशेषतः,
  $\Sat{M}{\lforall[x][!A(x)]}$.
  \olref[syn][sem]{prop:quant-terms} नुसार $\Sat{M}{!A(t)}$. परिणामी,
  $\Struct{M}$ हे $\sFmla{\True}{!A(t)}$ पूर्ण करते.

\item $\sFmla{\False}{\lforall[x][!A(x)]} \in \Gamma$ ला
  $\TRule{\False}{\lforall}$ लावून शाखेचा विस्तार केला आहे. यामुळे शाखेवर
  नवीन !!{signed formula}~$\sFmla{\False}{!A(a)}$ मिळते, जेथे $a$ हे
  $\Gamma$ मध्ये न येणारे !!a{constant} आहे. $\Gamma$ पूर्ततायोग्य असल्यामुळे
  $\Sat{M}{\Gamma}$ अशी एखादी $\Struct{M}$ असते; विशेषतः,
  $\Sat/{M}{\lforall[x][!A(x)]}$. आता
  $\Gamma \cup \{\sFmla{\False}{!A(a)}\}$ पूर्ततायोग्य आहे हे दाखवायचे आहे.
  त्यासाठी पुढीलप्रमाणे योग्य~$\Struct{M'}$ परिभाषित करू.

  \olref[syn][ass]{prop:sat-quant} नुसार,
  $\Sat/{M}{\lforall[x][!A(x)]}$ तेव्हा आणि केवळ तेव्हाच, जेव्हा काही $s$
  साठी $\Sat/{M}{!A(x)}[s]$. आता $\Struct{M'}$ ही $\Struct{M}$ सारखीच
  असू द्या, फक्त $\Assign{a}{M'} = s(x)$. $a$ हे $\Gamma$ मध्ये येत
  नसल्यामुळे, \olref[syn][ext]{cor:extensionality-sent} नुसार कोणत्याही
  $\sFmla{\True}{!C} \in \Gamma$ साठी $\Sat{M'}{!C}$, आणि कोणत्याही
  $\sFmla{\False}{!C} \in \Gamma$ साठी $\Sat/{M'}{!C}$.

  \olref[syn][ext]{prop:extensionality} नुसार $\Sat/{M'}{!A(x)}[s]$.
  \olref[syn][ext]{prop:ext-formulas} नुसार $\Sat/{M'}{!A(a)}[s]$.
  $!A(a)$ हे !!a{sentence} असल्यामुळे,
  \olref[syn][ass]{prop:sentence-sat-true} नुसार $\Sat/{M'}{!A(a)}$;
  म्हणजे $\Struct{M'}$ हे $\sFmla{\False}{!A(a)}$ पूर्ण करते.
\item $\sFmla{\True}{\lexists[x][!B(x)]} \in \Gamma$ ला
  $\TRule{\True}{\lexists}$ लावून शाखेचा विस्तार केला आहे: सराव.
\item $\sFmla{\False}{\lexists[x][!B(x)]} \in \Gamma$ ला
  $\TRule{\False}{\lexists}$ लावून शाखेचा विस्तार केला आहे: सराव.}{}
\end{enumerate}
आता शाखा विभागणारी संभाव्य अनुमाने पाहू.
\begin{enumerate}
\item $\sFmla{\False}{!B \land !C} \in \Gamma$ ला
  $\TRule{\False}{\land}$ लावून शाखेचा विस्तार केला आहे. यामुळे दोन शाखा
  मिळतात: डावीकडील शाखा $\sFmla{\False}{!B}$ मधून आणि उजवीकडील शाखा
  $\sFmla{\False}{!C}$ मधून पुढे जाते. समजा
  $\iftag{FOL}{\Sat{M}}{\pSat{v}}{\Gamma}$; विशेषतः,
  $\iftag{FOL}{\Sat/{M}}{\pSat/{v}}{!B \land !C}$. मग एकतर
  $\iftag{FOL}{\Sat/{M}}{\pSat/{v}}{!B}$ किंवा
  $\iftag{FOL}{\Sat/{M}}{\pSat/{v}}{!C}$. पहिल्या प्रकरणात
  $\iftag{FOL}{\Struct{M}}{\pAssign{v}}$ हे
  $\sFmla{\False}{!B}$ पूर्ण करते; म्हणजे
  $\iftag{FOL}{\Struct{M}}{\pAssign{v}}$ डावीकडील शाखेवरील सूत्रे पूर्ण करते.
  दुसऱ्या प्रकरणात $\iftag{FOL}{\Struct{M}}{\pAssign{v}}$ हे
  $\sFmla{\False}{!C}$ पूर्ण करते; म्हणजे
  $\iftag{FOL}{\Struct{M}}{\pAssign{v}}$ उजवीकडील शाखेवरील सूत्रे पूर्ण करते.
\item $\sFmla{\True}{!B \lor !C} \in \Gamma$ ला
  $\TRule{\True}{\lor}$ लावून शाखेचा विस्तार केला आहे: सराव.
\item $\sFmla{\True}{!B \lif !C} \in \Gamma$ ला
  $\TRule{\True}{\lif}$ लावून शाखेचा विस्तार केला आहे: सराव.
\item \Cut{} लावून शाखेचा विस्तार केला आहे. यामुळे दोन शाखा मिळतात: एका
  शाखेत $\sFmla{\True}{!B}$ आणि दुसऱ्या शाखेत $\sFmla{\False}{!B}$ असते.
  $\iftag{FOL}{\Sat{M}}{\pSat{v}}{\Gamma}$ आणि एकतर
  $\iftag{FOL}{\Sat{M}}{\pSat{v}}{!B}$ किंवा
  $\iftag{FOL}{\Sat/{M}}{\pSat/{v}}{!B}$ असल्यामुळे,
  $\iftag{FOL}{\Struct{M}}{\pAssign{v}}$ हे डावीकडील किंवा उजवीकडील शाखा
  पूर्ण करते.
\end{enumerate}
\end{proof}

\tagprob{FOL}
\begin{prob}
\olref[fol][tab][sou]{thm:tableau-soundness} ची सिद्धता पूर्ण करा.
\end{prob}
\tagendprob

\tagprob{notFOL}
\begin{prob}
\olref[pl][tab][sou]{thm:tableau-soundness} ची सिद्धता पूर्ण करा.
\end{prob}
\tagendprob

\begin{cor}
\ollabel{cor:weak-soundness}
जर $\Proves !A$, तर $!A$ हे \iftag{FOL}{वैध}{सर्वतः सत्य सूत्र} आहे.
\end{cor}

\begin{cor}
\ollabel{cor:entailment-soundness}
जर $\Gamma \Proves !A$, तर $\Gamma \Entails !A$.
\end{cor}

\begin{proof}
  जर $\Gamma \Proves !A$, तर काही $!B_1$, \dots, $!B_n \in
  \Gamma$ साठी $\{\sFmla{\False}{!A}, \sFmla{\True}{!B_1}, \dots,
  \sFmla{\True}{!B_n}\}$ चा बंद !!{tableau} असतो.
  \olref{thm:tableau-soundness} नुसार, प्रत्येक
  \iftag{FOL}{!!{structure}}{!!{valuation}}~$\iftag{FOL}{\Struct{M}}{\pAssign{v}}$
  काही $!B_i$ असत्य करते किंवा $!A$ सत्य करते. म्हणून
  $\iftag{FOL}{\Sat{M}}{\pSat{v}}{\Gamma}$ असेल, तर
  $\iftag{FOL}{\Sat{M}}{\pSat{v}}{!A}$ सुद्धा असेल.
\end{proof}

\begin{cor}
\ollabel{cor:consistency-soundness}
जर $\Gamma$ पूर्ततायोग्य असेल, तर तो सुसंगत आहे.
\end{cor}

\begin{proof}
प्रतिलोम विधान सिद्ध करू. $\Gamma$ सुसंगत नाही असे समजा. मग
$!B_1$, \dots, $!B_n \in \Gamma$ अशी वाक्ये आणि
$\{\sFmla{\True}{!B_1}, \dots, \sFmla{\True}{!B_n}\}$ चा बंद
!!{tableau} असतो. \olref{thm:tableau-soundness} नुसार, प्रत्येक
$i=1$, \dots,~$n$ साठी $\iftag{FOL}{\Sat{M}}{\pSat{v}}{!B_i}$ अशी
कोणतीही $\iftag{FOL}{\Struct{M}}{\pAssign{v}}$ नसते. पण मग $\Gamma$
पूर्ततायोग्य नसतो.
\end{proof}

\end{document}
